% 本书自编查阅表；全部公式与条件来自第三章已证明结果。
\begin{table}[H]
\centering\small\renewcommand{\arraystretch}{1.2}\setlength{\tabcolsep}{4pt}
\caption[Fourier 变换的运算规则]{Fourier 变换的运算规则。\(d\) 为正整数，\(x,\xi\in\mathbb R^d\)，统一采用指数 \(e^{-2\pi\mathrm{i}x\cdot\xi}\)；允许复值函数，实值情形包含在内。第三列列出各行的可积性与参数条件。}
\label{b2:tab:03-transform-rules}
\begin{tabularx}{\linewidth}{@{}>{\raggedright\arraybackslash}p{3.4cm}>{\raggedright\arraybackslash}p{6.1cm}>{\raggedright\arraybackslash}X@{}}
\toprule
空间中的操作 & 频率中的表达 & 适用条件 \\
\midrule
\(f(x-h)\) & \(e^{-2\pi\mathrm{i}h\cdot\xi}\cdot\widehat f(\xi)\) & \(f\in L^1(\mathbb R^d)\)，\(h\in\mathbb R^d\) \\
\(e^{2\pi\mathrm{i}\eta\cdot x}\cdot f(x)\) & \(\widehat f(\xi-\eta)\) & \(f\in L^1(\mathbb R^d)\)，\(\eta\in\mathbb R^d\) \\
\(f(Ax)\) & \(|\det A|^{-1}\cdot\widehat f(A^{-\mathsf T}\xi)\) & \(f\in L^1(\mathbb R^d)\)，\(A\in\mathbb R^{d\times d}\)，\(\det A\ne0\) \\
\(\partial_jf\) & \(2\pi\mathrm{i}\xi_j\widehat f(\xi)\) & \(1\leq j\leq d\)，\(f\in L^1(\mathbb R^d)\)，弱导数 \(\partial_jf\in L^1(\mathbb R^d)\) \\
\(x_jf\) & \(\dfrac{\mathrm{i}}{2\pi}\partial_{\xi_j}\widehat f\) & \(1\leq j\leq d\)，\(f,x_jf\in L^1(\mathbb R^d)\) \\
\(f*g\) & \(\widehat f\cdot\widehat g\) & \(f,g\in L^1(\mathbb R^d)\) \\
\(\begin{gathered}G_\varepsilon(x)=\\\varepsilon^{-d}e^{-\pi|x|^2/\varepsilon^2}\end{gathered}\) & \(e^{-\pi\varepsilon^2|\xi|^2}\) & \(\varepsilon>0\)，此核可积且积分为 \(1\) \\
\bottomrule
\end{tabularx}
\end{table}
